\begin{ouverture}[lewis]{5}{Électrons et stabilité}{Constitution et transformations}{Observer · Modéliser · Expérimenter · Expliquer}
\objectif{Lire une configuration électronique pour Z inférieur ou égal à 18.}
\objectif{Relier familles, électrons de valence et ions usuels.}
\objectif{Interpréter un schéma de Lewis fourni.}
\prerequis{Proportionnalité, conversions et lecture de graphiques.}
\end{ouverture}
\essentiel
\essSection{1}{Remplir les sous-couches}
Pour les atomes étudiés jusqu'à $Z=18$, les sous-couches se remplissent dans l'ordre 1s, 2s, 2p, 3s, 3p ; leurs capacités maximales sont respectivement 2, 2, 6, 2 et 6 électrons. La somme des exposants vaut le nombre d'électrons.
\begin{exemple}[Le soufre]Pour $Z=16$, la configuration est $1s^2\,2s^2\,2p^6\,3s^2\,3p^4$. La couche externe est celle de nombre 3 : elle contient 6 électrons de valence.\end{exemple}
La plus grande valeur du nombre de couche occupée donne la période. Pour les éléments des blocs s et p, les électrons de valence permettent de retrouver la colonne, avec le cas particulier de l'hélium, gaz noble à deux électrons.
\schema[corrige]{tableau-periodique}
\essSection{2}{Des entités plus stables}
Les gaz nobles ont une couche de valence saturée : deux électrons pour l'hélium, huit pour le néon et l'argon. Cette configuration aide à prévoir des ions monoatomiques usuels : le sodium perd un électron ; le magnésium en perd deux ; le fluor et le chlore en gagnent un.
\begin{definition}[Noms à connaître]$\ce{H+}$ : ion hydrogène ; $\ce{Na+}$ : sodium ; $\ce{K+}$ : potassium ; $\ce{Ca^{2+}}$ : calcium ; $\ce{Mg^{2+}}$ : magnésium ; $\ce{Cl-}$ : chlorure ; $\ce{F-}$ : fluorure.\end{definition}
Les membres d'une famille, dans une même colonne, ont des propriétés chimiques communes liées à leurs électrons de valence.
\essSection{3}{Lire les schémas de Lewis}
Un trait entre deux atomes représente un doublet liant, donc deux électrons partagés. Un doublet non liant appartient à un seul atome. Pour compter les électrons autour d'un atome, chaque doublet voisin compte pour deux, même s'il est partagé.
\schema[corrige]{lewis}
\begin{exemple}Dans le schéma fourni de l'eau, l'oxygène possède deux doublets liants et deux non liants : huit électrons l'entourent. Chaque hydrogène est entouré de deux électrons.\end{exemple}
Rompre une liaison demande de l'énergie. Une liaison plus difficile à rompre a une énergie de liaison plus élevée ; former des liaisons libère de l'énergie. En Seconde, on interprète les schémas fournis : aucune méthode générale de construction complexe n'est nécessaire.
\essSection{4}{Expérimenter et modéliser}
\begin{experience}[Confronter des modèles]Utiliser le tableau périodique et les schémas fournis pour compter les électrons de valence puis les doublets autour de chaque atome. Présenter un tableau de vérification pour \ce{H2}, \ce{H2O}, \ce{NH3} et \ce{CO2}. Les modèles de boules montrent l'enchaînement des atomes, mais pas tous les doublets.\end{experience}
\exercices
\rubrique{Prendre confiance}
\begin{blocLivre}{ex-5-1}\begin{exercice}[1]{Additionner les exposants}
La configuration d’un atome neutre est $1s^2\,2s^2\,2p^5$.
\schema{tableau-periodique}
\begin{questions}
\item Combien a-t-il d’électrons ?
\item Identifier l’élément avec le tableau fourni.
\end{questions}
\end{exercice}\end{blocLivre}
\begin{blocLivre}{ex-5-2}\begin{exercice}[1]{Électrons de valence}
Le magnésium a pour configuration $1s^2\,2s^2\,2p^6\,3s^2$.
\begin{questions}
\item Donner la période et le nombre d’électrons de valence.
\item Prévoir l’ion monoatomique usuel.
\end{questions}
\end{exercice}\end{blocLivre}
\begin{blocLivre}{ex-5-3}\begin{exercice}[1]{Nom et formule}
On donne les noms sodium, chlorure, calcium et fluorure pour des ions usuels.
\begin{questions}
\item Écrire les quatre formules ioniques.
\item Classer les ions en cations et anions.
\end{questions}
\end{exercice}\end{blocLivre}
\begin{blocLivre}{ex-5-4}\begin{exercice}[1]{Interpréter l’eau}
Utiliser le schéma de Lewis fourni pour \ce{H2O}.
\schema{lewis}
\begin{questions}
\item Compter les doublets liants et non liants de l’oxygène.
\item Vérifier le duet de chaque hydrogène.
\end{questions}
\end{exercice}\end{blocLivre}
\rubrique{S’entraîner}
\begin{blocLivre}{ex-5-5}\begin{exercice}[2]{Écrire une configuration}
Le silicium a $Z=14$.
\begin{questions}
\item Écrire sa configuration fondamentale.
\item Donner sa période et ses électrons de valence.
\end{questions}
\end{exercice}\end{blocLivre}
\begin{blocLivre}{ex-5-6}\begin{exercice}[2]{Deux membres d’une famille}
L’oxygène est $1s^2\,2s^2\,2p^4$ ; le soufre est $1s^2\,2s^2\,2p^6\,3s^2\,3p^4$.
\begin{questions}
\item Quel point commun de valence présentent-ils ?
\item Pourquoi sont-ils dans une même colonne ?
\end{questions}
\end{exercice}\end{blocLivre}
\begin{blocLivre}{ex-5-7}\begin{exercice}[2]{Un octet dans le dioxyde de carbone}
Le schéma de \ce{CO2} est fourni.
\schema{lewis}
\begin{questions}
\item Combien de doublets liants entourent le carbone ?
\item Vérifier l’octet du carbone et d’un oxygène.
\end{questions}
\end{exercice}\end{blocLivre}
\begin{blocLivre}{ex-5-8}\begin{exercice}[2]{Comparer des liaisons}
Pour deux liaisons isolées, on fournit des énergies de rupture de $3,0\times10^{-19}$ J et $5,0\times10^{-19}$ J.
\begin{questions}
\item Laquelle demande le plus d’énergie pour être rompue ?
\item Rompre une liaison libère-t-il de l’énergie ?
\end{questions}
\end{exercice}\end{blocLivre}
\rubrique{Raisonner et approfondir}
\begin{blocLivre}{ex-5-9}\begin{exercice}[3]{Un gaz noble particulier}
L’hélium possède $1s^2$ ; le néon $1s^2\,2s^2\,2p^6$.
\begin{questions}
\item Pourquoi sont-ils tous deux chimiquement stables ?
\item Faut-il huit électrons autour de l’hélium ?
\end{questions}
\end{exercice}\end{blocLivre}
\begin{blocLivre}{ex-5-10}\begin{exercice}[3]{Corriger un remplissage}
Une proposition pour le sodium est $1s^2\,2s^2\,2p^7$.
\begin{questions}
\item Identifier l’erreur.
\item Donner la configuration correcte et l’ion usuel.
\end{questions}
\end{exercice}\end{blocLivre}
\begin{blocLivre}{ex-5-11}\begin{exercice}[3]{La molécule d’ammoniac}
Le schéma fourni représente \ce{NH3}.
\schema{lewis}
\begin{questions}
\item Compter les électrons autour de l’azote.
\item Compter le nombre total de doublets figurés dans la molécule.
\end{questions}
\end{exercice}\end{blocLivre}
\begin{blocLivre}{ex-5-12}\begin{exercice}[3]{Atome ou ion ?}
Deux entités ont dix électrons. Leurs noyaux possèdent respectivement 10 et 12 protons.
\begin{questions}
\item Identifier l’atome neutre et le cation.
\item Que partagent leurs configurations électroniques ?
\end{questions}
\end{exercice}\end{blocLivre}

